\section{The electronics wall}

Chips are the part a lone factory cannot make, and the reason is twofold.

\subsection{Supply-chain depth}
The first foundation is supply-chain depth. Semiconductors sit at the end of the deepest
supply chain on Earth: more than 400 process steps, materials purified to nine-nines
purity, and a global monopoly on the extreme-ultraviolet lithography that patterns the
smallest features \cite{csis-semiconductor-supply-chain}. No single factory, however well
closed on structure, reproduces that pipeline from local rock.

\subsection{Embodied energy}
The second foundation is energy, and it is the one that turns a supply-chain argument into
a physical wall. The embodied energy of finished compute logic is on the order of
8{,}000~kWh/kg, against roughly 5~kWh/kg for structural metal smelted from regolith, a gap
of about three orders of magnitude \cite{williams-ayres-heller-2002,nagapurkar-das-2022}.
The gap is not an artifact of one study: an exergy analysis across twenty manufacturing
processes finds electricity per kilogram rising by orders of magnitude from conventional
metal shaping to vapor-phase semiconductor processing \cite{gutowski-2009}, and
process-level life-cycle assessments of CMOS logic across technology generations anchor the
finished-chip figure independently \cite{boyd-2012}. Table~\ref{tab:embodied} collects the
per-class figures with their sources and ranges, and Fig.~\ref{fig:embodied} shows the same
data across the seed factory's own subsystems. The measurement basis is what makes the
chip figure appear to move: a bare wafer is nearer 1{,}800~kWh/kg while a finished packaged
part runs to thousands \cite{murphy-2003}, so we fix the basis at the finished packaged
part throughout. Structural metals sit at the cheap end on independent datasheets
\cite{ice-coefficients,ashby-2012}, and in-situ regolith processing keeps them there off
Earth \cite{guerrero-zabel-2023}. Between the extremes lie the other electronics a factory
would rather not make: silicon photovoltaics near 40 to 120~kWh/kg \cite{peng-2013-pv-lca},
power electronics near 1{,}000 to 3{,}000~kWh/kg \cite{power-electronics-lca}, and sensors
near 2{,}000 to 8{,}000~kWh/kg \cite{sensor-embodied-2021}.

\begin{table}[!t]
  \centering
  \caption{Embodied electrical energy to manufacture one kilogram of a finished part, by
    class. Chips are about three orders of magnitude above structural metal.}
  \label{tab:embodied}
  \begin{tabular}{lcl}
    \toprule
    Material class & Embodied energy (kWh/kg) & Source \\
    \midrule
    Structural metal (smelted) & $\approx 5$    & \cite{ice-coefficients,ashby-2012} \\
    Silicon PV array           & 40-120         & \cite{peng-2013-pv-lca} \\
    Power electronics          & 1{,}000-3{,}000 & \cite{power-electronics-lca} \\
    Electronic sensors         & 2{,}000-8{,}000 & \cite{sensor-embodied-2021} \\
    Compute logic chips        & $\approx 8{,}000$ & \cite{williams-ayres-heller-2002,nagapurkar-das-2022,boyd-2012} \\
    \bottomrule
  \end{tabular}
\end{table}

\subsection{The power conditional}
These two foundations combine into a single conditional. A factory escapes the electronics
wall only if it is both highly self-sufficient \emph{and} swimming in power, and the
coupling is sharp. For a highly closed lunar seed at 97\% closure, letting it make its own
chips locally cuts the time to reach a target scale from about 29 years to about 17 years,
but only when it is fed on the order of 4~MW; at about 1~MW, making chips locally
backfires, and the factory runs out of electricity before it runs out of parts
\cite{nasa-cp-2255-1980}. Fig.~\ref{fig:crossover} shows the crossover: importing chips is
resupply-limited and stays near 29 years regardless of power, while making them locally is
energy-limited, beating the import time only above a threshold power and never completing
below it. Power is not a free knob. The energy cost of the chips is what decides whether
local production helps or hurts, and only a factory already swimming in power comes out
ahead.

\begin{figure}[!t]
  \centering
  \includegraphics[width=\columnwidth]{fig_embodied_energy.pdf}
  \caption{Embodied electrical energy to manufacture one kilogram of each seed-factory
    subsystem on-site (log scale). Chips and power electronics cost roughly 2{,}500 to
    8{,}000~kWh/kg, about three orders of magnitude more than the 3 to 7~kWh/kg of
    structural metal smelted from regolith, which is why electronics are the part a lone
    factory cannot afford to make.}
  \label{fig:embodied}
\end{figure}

\begin{figure}[!t]
  \centering
  \includegraphics[width=\columnwidth]{fig_chip_crossover.pdf}
  \caption{Time to reach the target factory scale versus available power for the 97\%-closed
    lunar seed. Importing chips (solid) is resupply-limited and stays near 29 years
    regardless of power; making chips locally (dashed) is energy-limited, reaching the
    target in about 17 years at 4~MW but diverging and never completing as power falls
    toward 1~MW.}
  \label{fig:crossover}
\end{figure}
